\section{OSPF NSSA Type-7 to Type-5 Translation}

\subsection{Objective}

The objective of this experiment was to validate OSPF NSSA behaviour, specifically Type-7 to Type-5 external route translation and external route containment.

The experiment checked whether an external route originated inside an NSSA area is first carried as a Type-7 LSA inside the NSSA, and then translated into a Type-5 LSA by the Area Border Router, or ABR, for the backbone area.

\subsection{NSSA Area in the Topology}

In this topology, Area 10 is configured as an NSSA area.

The relevant routers are:

\begin{table}[h]
\centering
\begin{tabular}{|l|l|}
\hline
\textbf{Router} & \textbf{Role} \\
\hline
\texttt{hpe-r3} & ABR between Area 0 and Area 10 \\
\hline
\texttt{hpe-r5} & Internal NSSA router \\
\hline
\texttt{hpe-r6} & Internal NSSA router and origin router for the test external route \\
\hline
\end{tabular}
\caption{Routers involved in the NSSA experiment}
\end{table}

Area 10 is configured as NSSA on \texttt{hpe-r3}, \texttt{hpe-r5}, and \texttt{hpe-r6}.

\subsection{Why NSSA is Important}

In normal OSPF, external routes are flooded as Type-5 LSAs.

However, an NSSA does not directly flood Type-5 LSAs inside the NSSA area. Instead, external routes originated inside the NSSA are carried as Type-7 LSAs. The ABR then translates the Type-7 LSA into a Type-5 LSA for the rest of the OSPF domain.

This keeps external route flooding controlled and prevents external LSAs from entering areas where they should not be directly flooded.

\subsection{Test External Route}

A controlled static blackhole route was added on \texttt{hpe-r6}.

\begin{verbatim}
172.16.66.0/24
\end{verbatim}

This was added as a BIRD static route:

\begin{verbatim}
172.16.66.0/24 blackhole
\end{verbatim}

This route was then exported into OSPF from \texttt{hpe-r6}, which is inside NSSA Area 10.

The route was intentionally configured as a blackhole route because the experiment only needed a clean routing advertisement, not real host connectivity.

\subsection{Route Visibility Result}

The route visibility result is shown below.

\begin{table}[h]
\centering
\begin{tabular}{|l|l|l|}
\hline
\textbf{Router} & \textbf{Route Present} & \textbf{Route Type} \\
\hline
\texttt{hpe-r6} & Yes & Static \\
\hline
\texttt{hpe-r5} & Yes & OSPF-E2 \\
\hline
\texttt{hpe-r3} & Yes & OSPF-E2 \\
\hline
\texttt{hpe-r1} & Yes & OSPF-E2 \\
\hline
\texttt{hpe-r2} & Yes & OSPF-E2 \\
\hline
\texttt{hpe-r4} & Yes & OSPF-E2 \\
\hline
\texttt{hpe-r8} & No & Not present \\
\hline
\end{tabular}
\caption{NSSA external route visibility}
\end{table}

This confirms that the test route was originated by \texttt{hpe-r6}, propagated through Area 10, translated by the ABR, and made visible to the backbone routers.

\subsection{Type-7 Evidence Inside NSSA}

The Type-7 LSADB output showed the test route inside the NSSA area.

\begin{verbatim}
hpe-r6:
0007  172.16.66.255   6.6.6.6

hpe-r5:
0007  172.16.66.255   6.6.6.6

hpe-r3:
0007  172.16.66.255   6.6.6.6
\end{verbatim}

This proves that the external route originated by \texttt{hpe-r6} was carried inside Area 10 as a Type-7 LSA.

\subsection{Type-5 Evidence Outside NSSA}

The Type-5 LSADB output showed the translated route outside the NSSA.

\begin{verbatim}
hpe-r3:
0005  172.16.66.255   3.3.3.3

hpe-r1:
0005  172.16.66.255   3.3.3.3

hpe-r2:
0005  172.16.66.255   3.3.3.3

hpe-r4:
0005  172.16.66.255   3.3.3.3
\end{verbatim}

This proves that ABR \texttt{hpe-r3}, whose router ID is \texttt{3.3.3.3}, translated the NSSA Type-7 LSA into a Type-5 LSA for the backbone area.

\subsection{Stub Area Containment}

Router \texttt{hpe-r8} is inside the stub area, Area 20.

When the route was checked on \texttt{hpe-r8}, BIRD reported:

\begin{verbatim}
Network not found
\end{verbatim}

This means \texttt{hpe-r8} did not receive the external OSPF route as a direct OSPF external route.

However, the Linux kernel still had a forwarding path using the default route:

\begin{verbatim}
172.16.66.1 via 10.0.78.2 dev eth0
\end{verbatim}

This is the expected containment behaviour. The external route itself was not flooded into the stub area, but the stub router could still forward traffic using the default route.

\subsection{Conclusion}

The NSSA experiment successfully demonstrated Type-7 to Type-5 translation and external route containment.

A static blackhole route, \texttt{172.16.66.0/24}, was originated inside NSSA Area 10 from \texttt{hpe-r6}. Inside the NSSA, the route appeared as a Type-7 LSA originated by router ID \texttt{6.6.6.6}. At the ABR, \texttt{hpe-r3}, the Type-7 LSA was translated into a Type-5 LSA originated by router ID \texttt{3.3.3.3}.

Backbone routers such as \texttt{hpe-r1}, \texttt{hpe-r2}, and \texttt{hpe-r4} received the translated Type-5 LSA. The stub-area router \texttt{hpe-r8} did not receive the external route directly, showing that external route flooding was contained.

This confirms that NSSA behaviour worked correctly in the lab topology.
